% Relativistic Modification Preamble
% Shared macros for all relativistic chapter modifications
%
% NOTE: This file requires \usepackage{tcolorbox} to be loaded
% in the main preamble (../../preamble.tex) for the custom
% environments defined at the end of this file.
%
% Usage: % Relativistic Modification Preamble
% Shared macros for all relativistic chapter modifications
%
% NOTE: This file requires \usepackage{tcolorbox} to be loaded
% in the main preamble (../../preamble.tex) for the custom
% environments defined at the end of this file.
%
% Usage: % Relativistic Modification Preamble
% Shared macros for all relativistic chapter modifications
%
% NOTE: This file requires \usepackage{tcolorbox} to be loaded
% in the main preamble (../../preamble.tex) for the custom
% environments defined at the end of this file.
%
% Usage: % Relativistic Modification Preamble
% Shared macros for all relativistic chapter modifications
%
% NOTE: This file requires \usepackage{tcolorbox} to be loaded
% in the main preamble (../../preamble.tex) for the custom
% environments defined at the end of this file.
%
% Usage: \input{rel_preamble} from any relativistic chapter
% or from rel_main.tex

% =====================================================================
% 4-Vector notation
% ---------------------------------------------------------------------
% Macros for writing 4-vectors and projection operators in a
% consistent style throughout all relativistic chapters.
% =====================================================================
\newcommand{\fourv}[1]{#1^{\mu}}
\newcommand{\fvel}{u^{\mu}}
\newcommand{\fveldown}{u_{\mu}}
\newcommand{\proj}{\Delta^{\mu\nu}}
\newcommand{\projdown}{\Delta_{\mu\nu}}

% =====================================================================
% Tensors
% ---------------------------------------------------------------------
% Energy-momentum tensor, Faraday tensor and its dual, metric tensors.
% These appear pervasively in the relativistic formulation of both
% hydrodynamic and hydromagnetic stability problems.
% =====================================================================
\newcommand{\emt}{T^{\mu\nu}}
\newcommand{\emtdown}{T_{\mu\nu}}
\newcommand{\farad}{F^{\mu\nu}}
\newcommand{\faraddown}{F_{\mu\nu}}
\newcommand{\dualF}{{}^{*}\!F^{\mu\nu}}
\newcommand{\metric}{g^{\mu\nu}}
\newcommand{\metricdown}{g_{\mu\nu}}
\newcommand{\mink}{\eta^{\mu\nu}}

% =====================================================================
% Covariant derivatives and Lie derivative
% ---------------------------------------------------------------------
% Covariant derivatives with different index labels for use in
% multi-index expressions without conflict.
% =====================================================================
\newcommand{\covd}{\nabla_{\mu}}
\newcommand{\covdu}{\nabla_{\nu}}
\newcommand{\covda}{\nabla_{\alpha}}
\newcommand{\Lie}{\mathcal{L}}

% =====================================================================
% Physical quantities
% ---------------------------------------------------------------------
% Core thermodynamic and kinematic quantities that appear in the
% relativistic equations of motion.  The enthalpy density
% w = epsilon + p plays the role that rho alone plays in the
% Newtonian theory.
% =====================================================================
\newcommand{\enthalpy}{w}           % w = epsilon + p
\newcommand{\edensity}{\varepsilon} % energy density
\newcommand{\rdensity}{\rho_0}     % rest-mass density
\newcommand{\Lf}{\gamma}            % Lorentz factor
\newcommand{\cs}{c_{\mathrm{s}}}    % sound speed
\newcommand{\vA}{v_{\mathrm{A}}}    % Alfven speed
\newcommand{\vf}{v_{\mathrm{f}}}    % fast magnetosonic speed
\newcommand{\vs}{v_{\mathrm{s}}}    % slow magnetosonic speed
\newcommand{\bfour}{b^{\mu}}        % magnetic 4-vector
\newcommand{\bsq}{b^2}             % b^mu b_mu (magnetic pressure term)

% =====================================================================
% Dissipation (Israel-Stewart causal theory)
% ---------------------------------------------------------------------
% The Israel-Stewart formalism replaces the acausal Navier-Stokes
% constitutive relations with hyperbolic relaxation equations.
% Each dissipative flux (bulk pressure Pi, shear tensor pi^{mu nu},
% heat flux q^mu) has its own relaxation time (tau_Pi, tau_pi, tau_q).
% =====================================================================
\newcommand{\bulkP}{\Pi}            % bulk viscous pressure
\newcommand{\shearten}{\pi^{\mu\nu}} % shear stress tensor
\newcommand{\heatflux}{q^{\mu}}     % heat flux 4-vector
\newcommand{\taupi}{\tau_{\pi}}     % shear relaxation time
\newcommand{\tauq}{\tau_{q}}        % heat relaxation time
\newcommand{\tauPi}{\tau_{\Pi}}     % bulk relaxation time
\newcommand{\bulkvisc}{\zeta}       % bulk viscosity coefficient
\newcommand{\shearvisc}{\eta_{\mathrm{s}}} % shear viscosity coefficient
\newcommand{\thermcond}{\kappa}     % thermal conductivity

% =====================================================================
% Dimensionless numbers (relativistic generalisations)
% ---------------------------------------------------------------------
% Relativistic analogues of the classical dimensionless parameters
% that govern stability thresholds.  Subscript "rel" distinguishes
% them from the Newtonian versions defined in the original chapters.
% =====================================================================
\newcommand{\Rarel}{\mathrm{Ra}_{\mathrm{rel}}}   % relativistic Rayleigh number
\newcommand{\Tarel}{\mathrm{Ta}_{\mathrm{rel}}}   % relativistic Taylor number
\newcommand{\Qrel}{Q_{\mathrm{rel}}}              % relativistic Chandrasekhar Q
\newcommand{\Mach}{\mathcal{M}}                    % Mach number

% =====================================================================
% Operators
% ---------------------------------------------------------------------
% The d'Alembertian (wave operator) replaces the Laplacian in the
% relativistic wave equations.  Two aliases are provided for
% author convenience.
% =====================================================================
\newcommand{\dAlembert}{\Box}       % d'Alembertian
\newcommand{\Dalembertian}{\Box}    % alias
\newcommand{\order}[1]{\mathcal{O}(#1)} % order-of-magnitude symbol

% =====================================================================
% Relativistic correction helpers
% ---------------------------------------------------------------------
% \relcorr{v^2/c^2} expands to [1 + O(v^2/c^2)], useful for
% writing correction factors.  \NRlimit produces an arrow
% labelled "v/c -> 0" for taking the non-relativistic limit.
% =====================================================================
\newcommand{\relcorr}[1]{\left[1 + \order{#1}\right]} % correction factor
\newcommand{\NRlimit}{\xrightarrow{v/c \to 0}}        % non-relativistic limit arrow

% =====================================================================
% Custom environments (require tcolorbox package)
% ---------------------------------------------------------------------
% causalitycheck  -- yellow box for verifying that perturbation
%                    equations preserve causality (signal speeds <= c).
% relcorrection   -- blue box highlighting where and how the
%                    relativistic result departs from the Newtonian one.
% =====================================================================
\newenvironment{causalitycheck}{%
  \begin{tcolorbox}[colback=yellow!5, colframe=yellow!50!black, title=Causality Verification]
}{%
  \end{tcolorbox}
}

\newenvironment{relcorrection}{%
  \begin{tcolorbox}[colback=blue!5, colframe=blue!50!black, title=Relativistic Correction]
}{%
  \end{tcolorbox}
}
 from any relativistic chapter
% or from rel_main.tex

% =====================================================================
% 4-Vector notation
% ---------------------------------------------------------------------
% Macros for writing 4-vectors and projection operators in a
% consistent style throughout all relativistic chapters.
% =====================================================================
\newcommand{\fourv}[1]{#1^{\mu}}
\newcommand{\fvel}{u^{\mu}}
\newcommand{\fveldown}{u_{\mu}}
\newcommand{\proj}{\Delta^{\mu\nu}}
\newcommand{\projdown}{\Delta_{\mu\nu}}

% =====================================================================
% Tensors
% ---------------------------------------------------------------------
% Energy-momentum tensor, Faraday tensor and its dual, metric tensors.
% These appear pervasively in the relativistic formulation of both
% hydrodynamic and hydromagnetic stability problems.
% =====================================================================
\newcommand{\emt}{T^{\mu\nu}}
\newcommand{\emtdown}{T_{\mu\nu}}
\newcommand{\farad}{F^{\mu\nu}}
\newcommand{\faraddown}{F_{\mu\nu}}
\newcommand{\dualF}{{}^{*}\!F^{\mu\nu}}
\newcommand{\metric}{g^{\mu\nu}}
\newcommand{\metricdown}{g_{\mu\nu}}
\newcommand{\mink}{\eta^{\mu\nu}}

% =====================================================================
% Covariant derivatives and Lie derivative
% ---------------------------------------------------------------------
% Covariant derivatives with different index labels for use in
% multi-index expressions without conflict.
% =====================================================================
\newcommand{\covd}{\nabla_{\mu}}
\newcommand{\covdu}{\nabla_{\nu}}
\newcommand{\covda}{\nabla_{\alpha}}
\newcommand{\Lie}{\mathcal{L}}

% =====================================================================
% Physical quantities
% ---------------------------------------------------------------------
% Core thermodynamic and kinematic quantities that appear in the
% relativistic equations of motion.  The enthalpy density
% w = epsilon + p plays the role that rho alone plays in the
% Newtonian theory.
% =====================================================================
\newcommand{\enthalpy}{w}           % w = epsilon + p
\newcommand{\edensity}{\varepsilon} % energy density
\newcommand{\rdensity}{\rho_0}     % rest-mass density
\newcommand{\Lf}{\gamma}            % Lorentz factor
\newcommand{\cs}{c_{\mathrm{s}}}    % sound speed
\newcommand{\vA}{v_{\mathrm{A}}}    % Alfven speed
\newcommand{\vf}{v_{\mathrm{f}}}    % fast magnetosonic speed
\newcommand{\vs}{v_{\mathrm{s}}}    % slow magnetosonic speed
\newcommand{\bfour}{b^{\mu}}        % magnetic 4-vector
\newcommand{\bsq}{b^2}             % b^mu b_mu (magnetic pressure term)

% =====================================================================
% Dissipation (Israel-Stewart causal theory)
% ---------------------------------------------------------------------
% The Israel-Stewart formalism replaces the acausal Navier-Stokes
% constitutive relations with hyperbolic relaxation equations.
% Each dissipative flux (bulk pressure Pi, shear tensor pi^{mu nu},
% heat flux q^mu) has its own relaxation time (tau_Pi, tau_pi, tau_q).
% =====================================================================
\newcommand{\bulkP}{\Pi}            % bulk viscous pressure
\newcommand{\shearten}{\pi^{\mu\nu}} % shear stress tensor
\newcommand{\heatflux}{q^{\mu}}     % heat flux 4-vector
\newcommand{\taupi}{\tau_{\pi}}     % shear relaxation time
\newcommand{\tauq}{\tau_{q}}        % heat relaxation time
\newcommand{\tauPi}{\tau_{\Pi}}     % bulk relaxation time
\newcommand{\bulkvisc}{\zeta}       % bulk viscosity coefficient
\newcommand{\shearvisc}{\eta_{\mathrm{s}}} % shear viscosity coefficient
\newcommand{\thermcond}{\kappa}     % thermal conductivity

% =====================================================================
% Dimensionless numbers (relativistic generalisations)
% ---------------------------------------------------------------------
% Relativistic analogues of the classical dimensionless parameters
% that govern stability thresholds.  Subscript "rel" distinguishes
% them from the Newtonian versions defined in the original chapters.
% =====================================================================
\newcommand{\Rarel}{\mathrm{Ra}_{\mathrm{rel}}}   % relativistic Rayleigh number
\newcommand{\Tarel}{\mathrm{Ta}_{\mathrm{rel}}}   % relativistic Taylor number
\newcommand{\Qrel}{Q_{\mathrm{rel}}}              % relativistic Chandrasekhar Q
\newcommand{\Mach}{\mathcal{M}}                    % Mach number

% =====================================================================
% Operators
% ---------------------------------------------------------------------
% The d'Alembertian (wave operator) replaces the Laplacian in the
% relativistic wave equations.  Two aliases are provided for
% author convenience.
% =====================================================================
\newcommand{\dAlembert}{\Box}       % d'Alembertian
\newcommand{\Dalembertian}{\Box}    % alias
\newcommand{\order}[1]{\mathcal{O}(#1)} % order-of-magnitude symbol

% =====================================================================
% Relativistic correction helpers
% ---------------------------------------------------------------------
% \relcorr{v^2/c^2} expands to [1 + O(v^2/c^2)], useful for
% writing correction factors.  \NRlimit produces an arrow
% labelled "v/c -> 0" for taking the non-relativistic limit.
% =====================================================================
\newcommand{\relcorr}[1]{\left[1 + \order{#1}\right]} % correction factor
\newcommand{\NRlimit}{\xrightarrow{v/c \to 0}}        % non-relativistic limit arrow

% =====================================================================
% Custom environments (require tcolorbox package)
% ---------------------------------------------------------------------
% causalitycheck  -- yellow box for verifying that perturbation
%                    equations preserve causality (signal speeds <= c).
% relcorrection   -- blue box highlighting where and how the
%                    relativistic result departs from the Newtonian one.
% =====================================================================
\newenvironment{causalitycheck}{%
  \begin{tcolorbox}[colback=yellow!5, colframe=yellow!50!black, title=Causality Verification]
}{%
  \end{tcolorbox}
}

\newenvironment{relcorrection}{%
  \begin{tcolorbox}[colback=blue!5, colframe=blue!50!black, title=Relativistic Correction]
}{%
  \end{tcolorbox}
}
 from any relativistic chapter
% or from rel_main.tex

% =====================================================================
% 4-Vector notation
% ---------------------------------------------------------------------
% Macros for writing 4-vectors and projection operators in a
% consistent style throughout all relativistic chapters.
% =====================================================================
\newcommand{\fourv}[1]{#1^{\mu}}
\newcommand{\fvel}{u^{\mu}}
\newcommand{\fveldown}{u_{\mu}}
\newcommand{\proj}{\Delta^{\mu\nu}}
\newcommand{\projdown}{\Delta_{\mu\nu}}

% =====================================================================
% Tensors
% ---------------------------------------------------------------------
% Energy-momentum tensor, Faraday tensor and its dual, metric tensors.
% These appear pervasively in the relativistic formulation of both
% hydrodynamic and hydromagnetic stability problems.
% =====================================================================
\newcommand{\emt}{T^{\mu\nu}}
\newcommand{\emtdown}{T_{\mu\nu}}
\newcommand{\farad}{F^{\mu\nu}}
\newcommand{\faraddown}{F_{\mu\nu}}
\newcommand{\dualF}{{}^{*}\!F^{\mu\nu}}
\newcommand{\metric}{g^{\mu\nu}}
\newcommand{\metricdown}{g_{\mu\nu}}
\newcommand{\mink}{\eta^{\mu\nu}}

% =====================================================================
% Covariant derivatives and Lie derivative
% ---------------------------------------------------------------------
% Covariant derivatives with different index labels for use in
% multi-index expressions without conflict.
% =====================================================================
\newcommand{\covd}{\nabla_{\mu}}
\newcommand{\covdu}{\nabla_{\nu}}
\newcommand{\covda}{\nabla_{\alpha}}
\newcommand{\Lie}{\mathcal{L}}

% =====================================================================
% Physical quantities
% ---------------------------------------------------------------------
% Core thermodynamic and kinematic quantities that appear in the
% relativistic equations of motion.  The enthalpy density
% w = epsilon + p plays the role that rho alone plays in the
% Newtonian theory.
% =====================================================================
\newcommand{\enthalpy}{w}           % w = epsilon + p
\newcommand{\edensity}{\varepsilon} % energy density
\newcommand{\rdensity}{\rho_0}     % rest-mass density
\newcommand{\Lf}{\gamma}            % Lorentz factor
\newcommand{\cs}{c_{\mathrm{s}}}    % sound speed
\newcommand{\vA}{v_{\mathrm{A}}}    % Alfven speed
\newcommand{\vf}{v_{\mathrm{f}}}    % fast magnetosonic speed
\newcommand{\vs}{v_{\mathrm{s}}}    % slow magnetosonic speed
\newcommand{\bfour}{b^{\mu}}        % magnetic 4-vector
\newcommand{\bsq}{b^2}             % b^mu b_mu (magnetic pressure term)

% =====================================================================
% Dissipation (Israel-Stewart causal theory)
% ---------------------------------------------------------------------
% The Israel-Stewart formalism replaces the acausal Navier-Stokes
% constitutive relations with hyperbolic relaxation equations.
% Each dissipative flux (bulk pressure Pi, shear tensor pi^{mu nu},
% heat flux q^mu) has its own relaxation time (tau_Pi, tau_pi, tau_q).
% =====================================================================
\newcommand{\bulkP}{\Pi}            % bulk viscous pressure
\newcommand{\shearten}{\pi^{\mu\nu}} % shear stress tensor
\newcommand{\heatflux}{q^{\mu}}     % heat flux 4-vector
\newcommand{\taupi}{\tau_{\pi}}     % shear relaxation time
\newcommand{\tauq}{\tau_{q}}        % heat relaxation time
\newcommand{\tauPi}{\tau_{\Pi}}     % bulk relaxation time
\newcommand{\bulkvisc}{\zeta}       % bulk viscosity coefficient
\newcommand{\shearvisc}{\eta_{\mathrm{s}}} % shear viscosity coefficient
\newcommand{\thermcond}{\kappa}     % thermal conductivity

% =====================================================================
% Dimensionless numbers (relativistic generalisations)
% ---------------------------------------------------------------------
% Relativistic analogues of the classical dimensionless parameters
% that govern stability thresholds.  Subscript "rel" distinguishes
% them from the Newtonian versions defined in the original chapters.
% =====================================================================
\newcommand{\Rarel}{\mathrm{Ra}_{\mathrm{rel}}}   % relativistic Rayleigh number
\newcommand{\Tarel}{\mathrm{Ta}_{\mathrm{rel}}}   % relativistic Taylor number
\newcommand{\Qrel}{Q_{\mathrm{rel}}}              % relativistic Chandrasekhar Q
\newcommand{\Mach}{\mathcal{M}}                    % Mach number

% =====================================================================
% Operators
% ---------------------------------------------------------------------
% The d'Alembertian (wave operator) replaces the Laplacian in the
% relativistic wave equations.  Two aliases are provided for
% author convenience.
% =====================================================================
\newcommand{\dAlembert}{\Box}       % d'Alembertian
\newcommand{\Dalembertian}{\Box}    % alias
\newcommand{\order}[1]{\mathcal{O}(#1)} % order-of-magnitude symbol

% =====================================================================
% Relativistic correction helpers
% ---------------------------------------------------------------------
% \relcorr{v^2/c^2} expands to [1 + O(v^2/c^2)], useful for
% writing correction factors.  \NRlimit produces an arrow
% labelled "v/c -> 0" for taking the non-relativistic limit.
% =====================================================================
\newcommand{\relcorr}[1]{\left[1 + \order{#1}\right]} % correction factor
\newcommand{\NRlimit}{\xrightarrow{v/c \to 0}}        % non-relativistic limit arrow

% =====================================================================
% Custom environments (require tcolorbox package)
% ---------------------------------------------------------------------
% causalitycheck  -- yellow box for verifying that perturbation
%                    equations preserve causality (signal speeds <= c).
% relcorrection   -- blue box highlighting where and how the
%                    relativistic result departs from the Newtonian one.
% =====================================================================
\newenvironment{causalitycheck}{%
  \begin{tcolorbox}[colback=yellow!5, colframe=yellow!50!black, title=Causality Verification]
}{%
  \end{tcolorbox}
}

\newenvironment{relcorrection}{%
  \begin{tcolorbox}[colback=blue!5, colframe=blue!50!black, title=Relativistic Correction]
}{%
  \end{tcolorbox}
}
 from any relativistic chapter
% or from rel_main.tex

% =====================================================================
% 4-Vector notation
% ---------------------------------------------------------------------
% Macros for writing 4-vectors and projection operators in a
% consistent style throughout all relativistic chapters.
% =====================================================================
\newcommand{\fourv}[1]{#1^{\mu}}
\newcommand{\fvel}{u^{\mu}}
\newcommand{\fveldown}{u_{\mu}}
\newcommand{\proj}{\Delta^{\mu\nu}}
\newcommand{\projdown}{\Delta_{\mu\nu}}

% =====================================================================
% Tensors
% ---------------------------------------------------------------------
% Energy-momentum tensor, Faraday tensor and its dual, metric tensors.
% These appear pervasively in the relativistic formulation of both
% hydrodynamic and hydromagnetic stability problems.
% =====================================================================
\newcommand{\emt}{T^{\mu\nu}}
\newcommand{\emtdown}{T_{\mu\nu}}
\newcommand{\farad}{F^{\mu\nu}}
\newcommand{\faraddown}{F_{\mu\nu}}
\newcommand{\dualF}{{}^{*}\!F^{\mu\nu}}
\newcommand{\metric}{g^{\mu\nu}}
\newcommand{\metricdown}{g_{\mu\nu}}
\newcommand{\mink}{\eta^{\mu\nu}}

% =====================================================================
% Covariant derivatives and Lie derivative
% ---------------------------------------------------------------------
% Covariant derivatives with different index labels for use in
% multi-index expressions without conflict.
% =====================================================================
\newcommand{\covd}{\nabla_{\mu}}
\newcommand{\covdu}{\nabla_{\nu}}
\newcommand{\covda}{\nabla_{\alpha}}
\newcommand{\Lie}{\mathcal{L}}

% =====================================================================
% Physical quantities
% ---------------------------------------------------------------------
% Core thermodynamic and kinematic quantities that appear in the
% relativistic equations of motion.  The enthalpy density
% w = epsilon + p plays the role that rho alone plays in the
% Newtonian theory.
% =====================================================================
\newcommand{\enthalpy}{w}           % w = epsilon + p
\newcommand{\edensity}{\varepsilon} % energy density
\newcommand{\rdensity}{\rho_0}     % rest-mass density
\newcommand{\Lf}{\gamma}            % Lorentz factor
\newcommand{\cs}{c_{\mathrm{s}}}    % sound speed
\newcommand{\vA}{v_{\mathrm{A}}}    % Alfven speed
\newcommand{\vf}{v_{\mathrm{f}}}    % fast magnetosonic speed
\newcommand{\vs}{v_{\mathrm{s}}}    % slow magnetosonic speed
\newcommand{\bfour}{b^{\mu}}        % magnetic 4-vector
\newcommand{\bsq}{b^2}             % b^mu b_mu (magnetic pressure term)

% =====================================================================
% Dissipation (Israel-Stewart causal theory)
% ---------------------------------------------------------------------
% The Israel-Stewart formalism replaces the acausal Navier-Stokes
% constitutive relations with hyperbolic relaxation equations.
% Each dissipative flux (bulk pressure Pi, shear tensor pi^{mu nu},
% heat flux q^mu) has its own relaxation time (tau_Pi, tau_pi, tau_q).
% =====================================================================
\newcommand{\bulkP}{\Pi}            % bulk viscous pressure
\newcommand{\shearten}{\pi^{\mu\nu}} % shear stress tensor
\newcommand{\heatflux}{q^{\mu}}     % heat flux 4-vector
\newcommand{\taupi}{\tau_{\pi}}     % shear relaxation time
\newcommand{\tauq}{\tau_{q}}        % heat relaxation time
\newcommand{\tauPi}{\tau_{\Pi}}     % bulk relaxation time
\newcommand{\bulkvisc}{\zeta}       % bulk viscosity coefficient
\newcommand{\shearvisc}{\eta_{\mathrm{s}}} % shear viscosity coefficient
\newcommand{\thermcond}{\kappa}     % thermal conductivity

% =====================================================================
% Dimensionless numbers (relativistic generalisations)
% ---------------------------------------------------------------------
% Relativistic analogues of the classical dimensionless parameters
% that govern stability thresholds.  Subscript "rel" distinguishes
% them from the Newtonian versions defined in the original chapters.
% =====================================================================
\newcommand{\Rarel}{\mathrm{Ra}_{\mathrm{rel}}}   % relativistic Rayleigh number
\newcommand{\Tarel}{\mathrm{Ta}_{\mathrm{rel}}}   % relativistic Taylor number
\newcommand{\Qrel}{Q_{\mathrm{rel}}}              % relativistic Chandrasekhar Q
\newcommand{\Mach}{\mathcal{M}}                    % Mach number

% =====================================================================
% Operators
% ---------------------------------------------------------------------
% The d'Alembertian (wave operator) replaces the Laplacian in the
% relativistic wave equations.  Two aliases are provided for
% author convenience.
% =====================================================================
\newcommand{\dAlembert}{\Box}       % d'Alembertian
\newcommand{\Dalembertian}{\Box}    % alias
\newcommand{\order}[1]{\mathcal{O}(#1)} % order-of-magnitude symbol

% =====================================================================
% Relativistic correction helpers
% ---------------------------------------------------------------------
% \relcorr{v^2/c^2} expands to [1 + O(v^2/c^2)], useful for
% writing correction factors.  \NRlimit produces an arrow
% labelled "v/c -> 0" for taking the non-relativistic limit.
% =====================================================================
\newcommand{\relcorr}[1]{\left[1 + \order{#1}\right]} % correction factor
\newcommand{\NRlimit}{\xrightarrow{v/c \to 0}}        % non-relativistic limit arrow

% =====================================================================
% Custom environments (require tcolorbox package)
% ---------------------------------------------------------------------
% causalitycheck  -- yellow box for verifying that perturbation
%                    equations preserve causality (signal speeds <= c).
% relcorrection   -- blue box highlighting where and how the
%                    relativistic result departs from the Newtonian one.
% =====================================================================
\newenvironment{causalitycheck}{%
  \begin{tcolorbox}[colback=yellow!5, colframe=yellow!50!black, title=Causality Verification]
}{%
  \end{tcolorbox}
}

\newenvironment{relcorrection}{%
  \begin{tcolorbox}[colback=blue!5, colframe=blue!50!black, title=Relativistic Correction]
}{%
  \end{tcolorbox}
}
